%!TEX root = ../main.tex
\section{Идея 4: Обучаем свой транслятор}

На этом я не остановился. Дальше я решил попробовать оставить в качестве главного планировщика Дейкстру и заменить транслятор траектории в интенцию своей нейросетью. До этого я использовал базовый транслятор из $\pi$-switch, но он был обучен универсально, а мне нужно было, чтобы интенция учитывала траекторию.

\subsection{Одноточечный транслятор}

Я решил обучить собственный специализированный транслятор $T_\theta(s_t, B(w_{\text{lookahead}})) \to z_{\text{cmd}}$, заточенный под перевод ближайшего вейпоинта в мягкую команду для $\pi^\ell$.

Я спроектировал модель с адаптивным гейтом $g_t = \sigma(W_g [e_s, e_w])$, 4 остаточными блоками и нормировкой на сферу $\mathbb{S}^{d-1}$. 

На лабиринте Medium это дало $83.3 \pm 5.0\%$ успеха, что чуть ниже, чем у базовой Дейкстры. Но когда мы перешли к более длинным лабиринтам, вскрылась фундаментальная проблема — \enquote{близорукость} транслятора: он точно так же, как и транслятор бейзлайна, не учитывал дальнейшую траекторию, а потому не мог обеспечить плавные повороты и своевременное замедление агента.

Это приводит нас к следующей, наиболее сильной и перспективной архитектуре.

\subsection{Трансформер цепочки вейпоинтов}

Раз у нас и так есть траектория из вейпоинтов и мы хотим использовать её для формирования интенций, то тут само собой просится механизм внимания на последовательность вейпоинтов $\mathcal{S} = [B(w_1), \dots, B(w_K), B(g)]$.

\begin{figure*}[t]
    \centering
    %!TEX root = ../main.tex
\begin{tikzpicture}[>=stealth, font=\small, scale=0.88, every node/.style={transform shape}]
    % Токены цепочки вейпоинтов
    \node[draw=blue!70, fill=blue!10, rounded corners=3pt, minimum width=1.6cm, align=center] (w1) at (0, 3.2) {$B(w_1)$\\$+ \cos \theta_1$};
    \node[draw=blue!70, fill=blue!10, rounded corners=3pt, minimum width=1.6cm, align=center] (w2) at (2.0, 3.2) {$B(w_2)$\\$+ \cos \theta_2$};
    \node[draw=blue!70, fill=blue!10, rounded corners=3pt, minimum width=1.6cm, align=center] (w3) at (4.0, 3.2) {$B(w_3)$\\$+ \cos \theta_3$};
    \node[draw=blue!70, fill=blue!10, rounded corners=3pt, minimum width=1.6cm, align=center] (w4) at (6.0, 3.2) {$B(w_4)$\\$+ \cos \theta_4$};
    \node[draw=orange!80, fill=orange!15, rounded corners=3pt, minimum width=1.6cm, align=center] (wg) at (8.4, 3.2) {$B(g)$\\(Финиш)};

    % Фигурная скобка: 4 локальных + 1 глобальный
    \draw[decorate, decoration={brace, amplitude=5pt}, thick] (7.0, 4.0) -- (-0.8, 4.0) node[midway, above=6pt, font=\scriptsize] {4 локальных вейпоинта};
    \draw[decorate, decoration={brace, amplitude=5pt}, thick] (9.4, 4.0) -- (7.4, 4.0) node[midway, above=6pt, font=\scriptsize] {Глобальная цель};

    % Блок Self-Attention
    \node[draw=blue!80, fill=blue!20, rounded corners=4pt, minimum width=10.2cm, minimum height=0.8cm, align=center] (self_attn) at (4.2, 1.8) {Trajectory Self-Attention (моделирование формы коридора)};

    % State Query
    \node[draw=green!70!black, fill=green!15, rounded corners=4pt, minimum width=2.2cm, minimum height=0.8cm, align=center] (obs) at (-1.5, 0.4) {Состояние $s_t$\\$\mathbb{R}^{29}$};

    % Cross-Attention с ALiBi
    \node[draw=purple!80, fill=purple!20, rounded corners=5pt, minimum width=8.5cm, minimum height=1.1cm, align=center] (cross_attn) at (5.0, 0.4) {\textbf{Multi-Head Cross-Attention со штрафом ALiBi}\\$\text{Attn}(Q, K) = \text{Softmax}\left(\frac{Q K^\top}{\sqrt{d}} - \lambda \cdot m_h \cdot [0, 1, 2, \dots]\right)$};

    % Проекция
    \node[draw=red!70!black, fill=red!15, rounded corners=4pt, minimum width=2.2cm, minimum height=0.8cm, align=center] (out) at (11.0, 0.4) {\textbf{Интенция}\\$z_{\text{cmd}} \to \pi^\ell$};

    \draw[->, thick] (w1) -- (0, 2.2);
    \draw[->, thick] (w2) -- (2.0, 2.2);
    \draw[->, thick] (w3) -- (4.0, 2.2);
    \draw[->, thick] (w4) -- (6.0, 2.2);
    \draw[->, thick] (wg) -- (8.4, 2.2);
    \draw[->, thick] (self_attn) -- (4.2, 0.95);
    \draw[->, thick] (obs) -- (cross_attn);
    \draw[->, thick] (cross_attn) -- (out);
\end{tikzpicture}

    \caption{Архитектура нашего лучшего метода Enhanced Sequence Waypoint Attention Transformer с затуханием ALiBi, токенами кривизны и окном $4+1$.}
    \label{fig:seq_transformer}
\end{figure*}

\subsection{Функция потерь транслятора последовательностей}

Транслятор обучался на следующем дифференцируемом лоссе:
\begin{equation}
\label{eq:seq_loss}
\begin{split}
\mathcal{L}(\theta) = {} & \underbrace{(1 - \cos(\hat{z}_t, z^*_t)) + 0.1 \|\hat{z}_t - z^*_t\|^2}_{\mathcal{L}_{\text{BC}}} \\
& + 0.5 \cdot \underbrace{\|\pi^\ell(s, \hat{z}) - a^*\|^2}_{\mathcal{L}_{\text{Action}}} - 0.02 \cdot \underbrace{F(s, \hat{z})^\top \hat{z}}_{\mathcal{L}_{\text{Reach}}} \\
& - 0.05 \cdot \underbrace{\hat{z}^\top B(w_1)}_{\mathcal{L}_{\text{Local}}} + 0.2 \cdot \underbrace{\|\hat{z}_t - z^*_{\text{single}}\|^2}_{\mathcal{L}_{\text{Aux}}}
\end{split}
\end{equation}
где $\mathcal{L}_{\text{BC}}$ выравнивает вектор интенции с учителем на сфере; $\mathcal{L}_{\text{Action}}$ дифференцируется сквозь замороженную политику $\pi^\ell$ и штрафует за ошибки в крутящих моментах моторов робота; $\mathcal{L}_{\text{Reach}}$ максимизирует меру физической проходимости пространства FB; $\mathcal{L}_{\text{Local}}$ ориентирует вектор команды на ближайший вейпоинт коридора $w_1$, а вспомогательный регуляризатор $\mathcal{L}_{\text{Aux}}$ служит жёстким локальным якорем, предотвращая перенос внимания Cross-Attention на далёкие точки за стеной на начальных эпохах.

В процессе экспериментов и доработки архитектуры я ввёл несколько ключевых улучшений:
\begin{enumerate}[leftmargin=*,noitemsep]
    \item \textbf{Токены кривизны поворотов ($\cos \theta_k$)}: по формуле
    \begin{equation}
        \label{eq:curvature_angle}
        \cos \theta_k = \frac{(w_k - w_{k-1}) \cdot (w_{k+1} - w_k)}{\|w_k - w_{k-1}\|_2 \|w_{k+1} - w_k\|_2 + 10^{-6}}
    \end{equation}
    мы заранее вычисляем косинус угла поворота между следующими сегментами пути. Агент заранее видит, прямой ли коридор впереди или резкий разворот на 90 градусов, и плавно меняет угол захода.
    \item \textbf{Затухание внимания ALiBi ($-\lambda m_h \mathbf{d}$)}: в первой версии трансформер начал \enquote{подглядывать} на далёкие точки за стеной и пытался спрямлять путь. Штраф ALiBi за порядковый номер заставляет сеть отдавать 80\% внимания ближайшим 2--3 точкам, сохраняя ориентацию на финиш лишь в хвосте.
    \item \textbf{Окно внимания}: мы ограничили внимание 4 ближайшими точками + терминальным вектором $B(g)$, чтобы сеть не перегружалась шумом далёких петель лабиринта.
    \item \textbf{Вспомогательный лосс локальной привязки $\mathcal{L}_{\text{Aux}}$}: предотвращает утечку внимания к точкам за стеной на ранних эпохах.
\end{enumerate}

\subsection{Результаты}
В результате на AntMaze Medium метод показал результат $83.6 \pm 3.8\%$ SR, что находится на одном уровне со всеми методами на основе Дейкстры. 
